\documentclass[11pt]{article}
\usepackage[margin=25mm]{geometry}
\usepackage{url}
\begin{document}
\begin{center}
{\Large Response to Reviewer 1 Conditions}\\
Paper No. S-28-4-002
\end{center}

We thank the reviewer and editor for the constructive conditional-acceptance comments. We revised the manuscript conservatively and addressed Reviewer 1 Conditions 1--13 one by one. In particular, we weakened the framing from a privacy-preserving in-the-wild physiological sensing system to a browser-native, local-first, task-aligned physiological feature logging framework evaluated in an exploratory N=1 author-participant home/laboratory deployment.

\section*{Condition 1: Clarify the contribution}
\textbf{Response:} Thank you for the comment. We revised the manuscript to state that Camerala is not a new rPPG estimator, not a dedicated privacy-preserving algorithm, and not a general in-the-wild sensing system.\par
\textbf{Revision:} The title, abstract, introduction, contribution list, related work, discussion, and conclusion now present Camerala as a browser-native local-first implementation example.\par
\textbf{Location:} Title; Abstract; Sections 1--3; Discussion; Conclusion.\par
\textbf{Key revised text:} ``The contribution of Camerala is therefore not a new rPPG estimator or a dedicated privacy-preserving algorithm.''

\section*{Condition 2: Weaken privacy claims}
\textbf{Response:} Thank you for the comment. We removed the title phrase ``Privacy-Preserving'' and replaced the claim with local-first and privacy-conscious architectural wording.\par
\textbf{Revision:} The manuscript now says raw video is not transmitted or stored by Camerala under the tested implementation, but does not claim formal privacy preservation.\par
\textbf{Location:} Title; Abstract; Introduction; System Architecture.\par
\textbf{Key revised text:} ``This does not constitute a dedicated privacy-preserving algorithm; rather, it is a system-design choice that reduces the need to transmit or retain raw video.''

\section*{Condition 3: Reformulate research questions}
\textbf{Response:} Thank you for the comment. We revised the RQs to focus on operational completion and descriptive measurement-quality observations.\par
\textbf{Revision:} RQ1 is limited to the tested laptop, browser, camera, and task configuration. RQ2 asks about descriptive differences in measurement-quality indicators and sensor-derived estimates.\par
\textbf{Location:} Section 1.1.\par
\textbf{Key revised text:} ``These criteria are intended only to describe operational completion under the tested configuration; they do not demonstrate general deployability across users, devices, or environments.''

\section*{Condition 4: Revise contributions}
\textbf{Response:} Thank you for the comment. We replaced stronger novelty claims with three conservative contributions.\par
\textbf{Revision:} The contributions are now browser-native local-first sensing architecture, task-aligned on-device sensing and local persistence, and an exploratory task-constrained home/laboratory deployment.\par
\textbf{Location:} Section 1.2.\par
\textbf{Key revised text:} ``In this paper, we focus on system integration and conservative deployment observations.''

\section*{Condition 5: Add reviewer-mentioned related works and reposition novelty}
\textbf{Response:} Thank you for the comment. We added and discussed the reviewer-mentioned works/systems.\par
\textbf{Revision:} We added Gupta and Etemad (2023), Ayesha, Qiao, and Zulkernine (2022), LabVanced rPPG, FacePhys-Demo, RhythmEdge, and Di Lernia et al. (2024). We explicitly state that browser-based, web-deployable, and edge-based rPPG systems already exist.\par
\textbf{Location:} Section 2.2; reference list.\par
\textbf{Key revised text:} ``Accordingly, local rPPG processing alone is not Camerala's novelty, and privacy-oriented edge processing alone is not Camerala's novelty.''

\section*{Condition 6: Avoid claims of physiological, cognitive, or motivational effects}
\textbf{Response:} Thank you for the comment. We revised Results and Discussion to use descriptive language.\par
\textbf{Revision:} The instruction blocks are now described as instructional task contexts rather than validated manipulations.\par
\textbf{Location:} Abstract; Methods; Results; Discussion; Limitations.\par
\textbf{Key revised text:} ``They are used only as different instructional task contexts for demonstrating synchronized logging across task conditions.''

\section*{Condition 7: Treat rPPG-derived outputs conservatively}
\textbf{Response:} Thank you for the comment. We removed revision-history wording and avoided interpreting bpm-style summaries as validated heart-rate changes.\par
\textbf{Revision:} The manuscript now states that no ECG or contact PPG reference was collected.\par
\textbf{Location:} Results, rPPG-oriented ROI estimates.\par
\textbf{Key revised text:} ``Because no ECG or contact PPG reference was collected, block-level bpm-style summaries are not interpreted as validated heart-rate changes in this paper.''

\section*{Condition 8: Revise figures and captions}
\textbf{Response:} Thank you for the comment. We removed overstrong labels and inferential markings from the revised figures.\par
\textbf{Revision:} Figure labels now use neutral wording such as head-motion estimates, reaction time, instructional task contexts, laboratory, and home. P-values and significance marks are not shown.\par
\textbf{Location:} Figures 2--4 and captions.\par
\textbf{Key revised text:} ``No p-values, significance marks, or psychological/physiological effect labels are shown.''

\section*{Condition 9: Add hardware/software reproducibility details}
\textbf{Response:} Thank you for the comment. We added a hardware/software table and distinguished logged facts from unavailable facts.\par
\textbf{Revision:} Available user-agent evidence, browser information, requested camera constraints, processing FPS reporting, and unlogged hardware/camera fields are now explicitly separated.\par
\textbf{Location:} Section 4.3, Table 2.\par
\textbf{Key revised text:} ``Unknown or unrecorded deployment-time details are explicitly marked rather than inferred from the current revision environment.''

\section*{Condition 10: Add environment details}
\textbf{Response:} Thank you for the comment. We added an environment and lighting table.\par
\textbf{Revision:} Known lux values and lighting descriptions are reported. Unrecoverable details are marked as not logged or not systematically recorded rather than inferred.\par
\textbf{Location:} Section 4.4, Table 3.\par
\textbf{Key revised text:} ``Table 3 separates reported facts from information that was not systematically recorded.''

\section*{Condition 11: Clarify FaceMesh confidence threshold}
\textbf{Response:} Thank you for the comment. We revised the confidence-threshold wording.\par
\textbf{Revision:} The threshold is now described as an operational exclusion threshold for FaceMesh tracking confidence, not as an rPPG-quality threshold.\par
\textbf{Location:} Section 4.7; Results.\par
\textbf{Key revised text:} ``This threshold was not independently validated as an rPPG-quality threshold; therefore, the resulting face-tracking availability is treated only as a FaceMesh/logging quality indicator.''

\section*{Condition 12: Expand limitations}
\textbf{Response:} Thank you for the comment. We expanded the limitations.\par
\textbf{Revision:} The limitations now include N=1, author-participant, no ECG/contact PPG, no reference physiological sensor, no multiple participants, no multiple devices, no multiple home environments, no controlled separation of environmental and task factors, and no comparison with existing systems.\par
\textbf{Location:} Section 6.5.\par
\textbf{Key revised text:} ``The study does not claim clinical or medical accuracy, general deployability, broad operational generality, cross-device stability, or validated cognitive, motivational, psychological, or physiological effects.''

\section*{Condition 13: Focus discussion on safeguards for future deployment}
\textbf{Response:} Thank you for the comment. We revised the Discussion to focus on design requirements and safeguards.\par
\textbf{Revision:} The Discussion now emphasizes local-first logging, measurement-quality indicators, interpretation under environmental heterogeneity, and safeguards for future Camerala deployments.\par
\textbf{Location:} Section 6.\par
\textbf{Key revised text:} ``Future deployments should incorporate clearer safeguards: documented hardware/browser/camera settings, explicit reporting of camera constraints and actual camera settings, pre-session lighting checks, face-size and tracking warnings, luminance and exposure-fluctuation flags, and predefined quality-gated analysis rules.''

\end{document}
